\section{Principios de LightRAG}

\subsection{Panorama general}

LightRAG es un artículo académico de EMNLP 2025, además de otro proyecto del autor de Paper2Slides. Es un sistema RAG basado en un \textbf{grafo de conocimiento}, con una capacidad de comprensión semántica más sólida que el RAG vectorial tradicional.

\begin{importantbox}{Características centrales de LightRAG}
\begin{itemize}
    \item Indexación basada en \textbf{Knowledge Graph} (en lugar de un índice puramente vectorial)
    \item \textbf{Recuperación en dos niveles}: Low-Level (local, entidades específicas) + High-Level (global, temas y conceptos)
    \item \textbf{Actualización incremental}: los documentos nuevos no requieren reconstruir todo el Graph; basta con ejecutar el paso de extracción y fusionar
\end{itemize}
\end{importantbox}

\subsection{Proceso de Indexing: construcción del grafo de conocimiento}

Para cada documento se ejecutan los siguientes pasos:

\begin{enumerate}
    \item \textbf{Chunking}: dividir el documento en fragmentos
    \item \textbf{Entity \& Relationship Extraction}: extraer entidades y relaciones de cada Chunk mediante un LLM
    \item \textbf{Profiling}: generar descripciones textuales de las entidades y relaciones mediante un LLM
    \item \textbf{Deduplication}: deduplicar y fusionar en todo el grafo de conocimiento
    \item \textbf{Embedding}: construir un Embedding Vector para la Key de todas las entidades y relaciones
\end{enumerate}

Estructura de datos del grafo de conocimiento:
\begin{itemize}
    \item \textbf{Entity}: nombre, tipo, descripción (generada por Profiling), ID del Chunk original
    \item \textbf{Edge}: Source, Target, palabras clave, descripción, ID del Chunk original
\end{itemize}

\begin{knowledgebox}{LightRAG frente a GraphRAG}
Una ventaja importante de LightRAG frente a GraphRAG es la \textbf{actualización incremental}: para un documento nuevo basta con ejecutar los mismos pasos de extracción y fusionar las entidades y los bordes. En cambio, GraphRAG necesita rehacer la detección de comunidades, lo cual resulta costoso.
\end{knowledgebox}

\subsection{Proceso de Retrieval: recuperación en dos niveles}

\begin{enumerate}
    \item A partir de la Query, el LLM genera \textbf{Low-Level Keywords} (entidades específicas) y \textbf{High-Level Keywords} (temas o conceptos)
    \item Se emparejan los nodos relevantes del grafo de conocimiento según la similitud de Embedding
    \item Se amplía el Subgraph mediante \textbf{vecinos a un salto}
    \item Se obtienen tres tipos de Context: entidades, relaciones, Chunk original
    \item Context + Query $\rightarrow$ el LLM genera la respuesta final
\end{enumerate}

\subsection{Resumen de la sección}

LightRAG logra, mediante la indexación por grafo de conocimiento y la recuperación en dos niveles, una capacidad de comprensión semántica más sólida que el RAG puramente vectorial, además de admitir actualización incremental.

